% !TEX program = xelatex
\documentclass[UTF8,a4paper,zihao=-4,fontset=fandol]{ctexart}
\usepackage[a4paper,left=2.5cm,right=2.5cm,top=2.3cm,bottom=2.3cm]{geometry}
\usepackage{booktabs,longtable,enumitem,hyperref}
\hypersetup{hidelinks}
\setlength{\parindent}{2em}
\setlength{\parskip}{0.35em}
\renewcommand{\arraystretch}{1.25}
\pagestyle{plain}
\title{AI工具使用详情}
\author{}
\date{}

\begin{document}
\maketitle

\section{所用AI工具名称、版本或型号}
竞赛过程中使用OpenAI Codex桌面应用提供的代码与文档辅助能力，并调用本地配置的文档解析、图表生成、LaTeX编译、PDF视觉检查等技能或插件。具体底层模型版本由平台运行环境管理；本文不虚构不可核验的模型编号。COPT求解器、Python科学计算库和LaTeX属于数值计算与排版工具，不作为生成式AI工具列入。

\section{具体使用目的和环节}
\begin{enumerate}[label=\arabic*.]
  \item 审题与任务拆解：协助逐条整理四问的数据流、公式、六类硬约束、输出模板和匿名提交规范。
  \item 代码实现与调试：协助编写Q1特征与时间验证、Q2/Q3 COPT模型、Q3蒙特卡洛模拟、Q4公开赛程解析、自动验收和图表脚本。
  \item 公式和结果核验：将官方题面解析文本与实现常量逐项对照，检查单位、归一化、静态/动态信息截面和目标函数。
  \item 论文整理：协助生成LaTeX初稿、参考文献、图表说明、复现附录和排版修复。
\end{enumerate}

\section{主要提示方式与使用过程说明}
主要提示以本地赛题文件、数据附件、全过程指导文件和当前工作树为依据，要求AI先读取题面与已有实现，再逐步完成Q1至Q4、COPT求解、稳健性、论文和提交包。关键约束包括：不得使用赛后变量预测Q1；不得改变固定train/test；Q2/Q3实质使用COPT并保留日志；变量必须有题面或事实依据；不得声称隐藏测试或全空间最优；每次修改后必须运行验证。使用过程采用“读取证据---提出或实现修改---运行代码---检查日志与结果---再写论文”的迭代方式，而非直接采纳一次性生成内容。

\section{AI输出的采纳、人工修改和核验}
参赛队对AI输出承担全部责任，并完成如下人工审查与可执行核验：
\begin{enumerate}[label=\arabic*.]
  \item 对全部Q1--Q4公式、变量来源、单位和输出表头与官方PDF及CSV模板逐项核对；删除或拒绝无事实依据的指标。
  \item 运行19项Python测试；运行14项最终交付验收，覆盖六份模板列序、行数、Q1单位、Q2六类硬约束、Q3 32个晋级名额与每日资源约束、Q4来源与赛事结构。
  \item Q2中将早期可加代理目标人工升级为完整受限候选$Z_2$线性化；COPT目标0.45816365037316586与独立评估器0.458163650373169对账，误差小于$4\times10^{-15}$。
  \item 保留Q2/Q3完整COPT日志，核对版本、变量/约束规模、整数状态、最优界、gap和约束违反量；Q3除静态、动态最大化模型外，还在同一1623个候选及每日约束下求解最小化模型，得到同口径可行目标区间$[4.7963,6.6522]$；明确Q2最优性仅限受限候选集合，Q3区间也仅限共同离散候选集。
  \item 用5000、10000、20000和50000次模拟检查Q3收敛；用10000组权重扰动检查Q2结论稳健性。
  \item 对Q4公开来源保留原始快照、URL、采集日期和SHA-256，不采纳无法从公开数据同口径识别的商业指标。
  \item 使用XeLaTeX/Biber编译论文并将每页渲染为PNG检查，修复图像裁切、表格重叠和字体缺字后再纳入提交材料；论文附录逐文件嵌入支撑材料中的全部35个Python与PowerShell源程序，并核对源码清单无缺失、无额外项。
\end{enumerate}

\section{未采纳或经修改的典型内容}
早期方案曾将Q2 COPT目标描述为与完整$Z_2$同方向的代理目标，后经人工审计发现其不足，改为对票务、转播、成本、旅行、风险、场馆启用成本及两类公平极差作完整线性化。早期Q4说明曾记录网络采集失败，数据获取成功后已按原始来源与哈希重新生成并更新文档。论文中没有采纳“全局最优”“隐藏测试表现优异”或“优化赛程全面优于真实赛程”等无法由证据支持的表述。

\end{document}
